\scanpage{109}

\[
(\deleted{\unsure{S}}\mathrm T_0\,2)\qquad P_0\text{ est orientable}\tag*{(24)}
\]

% slanted
\margin{Il n'est pas sûr que cette condition en (24) soit pas prématurée -- \unsure{s'il ne|si elle} \deleted{\unsure{vaudrait}} vaut pas mieux l'ignorer d'abord, et introduire $Z'$ comme commutant, non de $G_0^+$, mais de $G_0$ tout entier.}

Je désigne par $\mu_0$ l'ensemble de ses orientations

\[
\left\{\begin{aligned}
 &\mu_0=\operatorname{Or}(P_0),\\
 &G_0^+=\operatorname{Ker}\bigl(G_0\longrightarrow\unsure{\mathfrak S|\mathfrak G}_{\mu_0}\bigr)\\
 &\text{groupe des autom. \uline{directs} du polyèdre}\\
 &P_0
\end{aligned}\right.\tag*{(25)}
\]

ainsi $P_0$ est simplement transitif sur l'ens $\vec A_0$ des arêtes orientées de $P_0$.

On introduit

\[
\begin{aligned}
 Z'&\subset\mathfrak S_S\\
 \shortparallel&\\
 \text{sous-groupe de }\mathfrak S_S&\text{ commutant de }G_0^+.
\end{aligned}\tag*{(26)}
\]

% slanted
\margin{Il \unsure{vaut|est} mieux réserver l'axiome $\mathrm T_0\,2$ pour plus tard, et introduire $Z'$ comme le commutant de $G_0$ dans $\mathfrak S$ -- ce qui dispense plus tard d'introduire l'axiome-\unsure{remords} $\mathrm T_0\,6$ que $\deleted{\unsure{Z_0}}Z'$ est bien \unsure{ég.} à ce commutant !}

Comme $S$ est un espace homogène \deleted{\unsure{sous}}\unsure{dans|sous} $\deleted{\unsure{G_0}}G_0^+$, choisissant une ``origine''\ $s\in S$ \unsure{ds|dans} cet espace homogène, et introduisant le stabilisateur de $s$ dans $G_0^+$, on a un isom.

\[
Z'\simeq\operatorname{Norm}_{G_0^+}(G_{0s}^+)/G_{0s}^+\tag*{(27)}
\]

\deleted{Remarquons}\inserted{Notons} que \inserted{si $n$ est l'ordre des sommets de $P_0$,} le groupe $G_{0\unsure{s|2}}^+$ est isomorphe à $\mathbb Z/n\mathbb Z=\Lambda_n$ \deleted{\unsure{(un}} : l'isomorphisme
% archive sheet number: 108
